\documentclass[12pt, a4paper, oneside]{book}

% =====================================================
% PACKAGES
% =====================================================
\usepackage[utf8]{inputenc}
\usepackage[left=0.55cm, right=0.55cm, top=0.7cm, bottom=0.7cm, includeheadfoot]{geometry}
\usepackage{graphicx}
\usepackage{xcolor}
\usepackage{listings}
\usepackage{fancyhdr}
\usepackage{hyperref}
\usepackage{titlesec}
\usepackage{tocloft}
\usepackage{amsmath}
\usepackage{amssymb}
\usepackage{array}
\usepackage{booktabs}
\usepackage{tikz}
\usepackage{mdframed}
\usepackage{longtable}
\usepackage{multirow}
\usepackage{colortbl}
\usepackage{tabularx}
\usetikzlibrary{shapes.geometric, arrows.meta, positioning, fit, backgrounds, shadows, decorations.pathreplacing}

% =====================================================
% COLOR DEFINITIONS
% =====================================================
\definecolor{darkblue}{RGB}{25, 55, 109}
\definecolor{accentblue}{RGB}{70, 130, 180}
\definecolor{lightblue}{RGB}{230, 240, 250}
\definecolor{darkgreen}{RGB}{34, 139, 34}
\definecolor{gold}{RGB}{218, 165, 32}
\definecolor{darkred}{RGB}{139, 0, 0}
\definecolor{orange}{RGB}{255, 140, 0}
\definecolor{purple}{RGB}{128, 0, 128}
\definecolor{teal}{RGB}{0, 128, 128}

% Instruction type colors
\definecolor{rtypecolor}{RGB}{229, 57, 53}
\definecolor{itypecolor}{RGB}{30, 136, 229}
\definecolor{stypecolor}{RGB}{67, 160, 71}
\definecolor{btypecolor}{RGB}{251, 140, 0}
\definecolor{utypecolor}{RGB}{142, 36, 170}
\definecolor{jtypecolor}{RGB}{0, 137, 123}

% Field colors (for tables)
\definecolor{funct7bg}{RGB}{255, 205, 210}
\definecolor{rs2bg}{RGB}{200, 230, 201}
\definecolor{rs1bg}{RGB}{187, 222, 251}
\definecolor{funct3bg}{RGB}{255, 224, 178}
\definecolor{rdbg}{RGB}{225, 190, 231}
\definecolor{opcodebg}{RGB}{179, 229, 252}
\definecolor{immbg}{RGB}{255, 205, 210}

% =====================================================
% ASSEMBLY SYNTAX HIGHLIGHTING
% =====================================================
\lstdefinelanguage{riscv}{
  morekeywords={
    add, sub, sll, slt, sltu, xor, srl, sra, or, and,
    addi, slti, sltiu, xori, ori, andi, slli, srli, srai,
    lb, lh, lw, lbu, lhu,
    sb, sh, sw,
    beq, bne, blt, bge, bltu, bgeu,
    lui, auipc,
    jal, jalr,
    ecall, ebreak,
    x0, x1, x2, x3, x4, x5, x6, x7, x8, x9,
    x10, x11, x12, x13, x14, x15, x16, x17, x18, x19,
    x20, x21, x22, x23, x24, x25, x26, x27, x28, x29,
    x30, x31,
    zero, ra, sp, gp, tp, t0, t1, t2, s0, fp, s1,
    a0, a1, a2, a3, a4, a5, a6, a7,
    s2, s3, s4, s5, s6, s7, s8, s9, s10, s11,
    t3, t4, t5, t6
  },
  sensitive=true,
  morecomment=[l]{\#},
  morecomment=[l]{//},
  morestring=[b]"
}

\lstset{
  language=riscv,
  basicstyle=\ttfamily\small,
  keywordstyle=\color{purple}\bfseries,
  commentstyle=\color{darkgreen}\itshape,
  stringstyle=\color{orange},
  numberstyle=\color{gray}\tiny,
  identifierstyle=\color{darkblue},
  backgroundcolor=\color{lightblue},
  breaklines=true,
  breakatwhitespace=true,
  tabsize=2,
  showstringspaces=false,
  numbers=left,
  numbersep=5pt,
  framexleftmargin=10pt,
  framextopmargin=5pt,
  framexbottommargin=5pt,
  frame=tb,
  rulecolor=\color{accentblue},
  captionpos=b,
  lineskip=-1pt
}

% =====================================================
% TITLE FORMATTING
% =====================================================
\titleformat{\chapter}
  {\normalfont\huge\bfseries\color{darkblue}}
  {\color{gold}\thechapter.}
  {20pt}
  {}
  [\color{accentblue}\hrule]
\titlespacing*{\chapter}{0pt}{-30pt}{20pt}

\titleformat{\section}
  {\normalfont\Large\bfseries\color{accentblue}}
  {\color{gold}\thesection}
  {12pt}
  {}

\titleformat{\subsection}
  {\normalfont\large\bfseries\color{teal}}
  {\color{gold}\thesubsection}
  {10pt}
  {}

% =====================================================
% HEADER AND FOOTER
% =====================================================
\pagestyle{fancy}
\fancyhf{}
\fancyhead[L]{\color{darkblue}\textbf{YARP RV32I --- Design \& Verification Guide}}
\fancyhead[R]{\color{darkblue}\thepage}
\fancyfoot[C]{\color{gray}\small\today}
\renewcommand{\headrulewidth}{1.5pt}
\renewcommand{\headrule}{\color{accentblue}\hrule}

% =====================================================
% TABLE OF CONTENTS STYLING
% =====================================================
\renewcommand{\cftchapfont}{\color{darkblue}\bfseries}
\renewcommand{\cftsecfont}{\color{accentblue}}
\renewcommand{\cftdot}{\color{gold}.}

% =====================================================
% HYPERLINK COLORS
% =====================================================
\hypersetup{
  colorlinks=true,
  linkcolor=darkblue,
  urlcolor=accentblue,
  citecolor=darkgreen
}

% =====================================================
% COMPONENT DESCRIPTION BOX
% =====================================================
\newmdenv[
  linecolor=accentblue,
  linewidth=1.5pt,
  backgroundcolor=lightblue!60,
  roundcorner=6pt,
  skipabove=6pt,
  skipbelow=6pt,
  innertopmargin=8pt,
  innerbottommargin=8pt,
  innerleftmargin=10pt,
  innerrightmargin=10pt
]{infobox}

% Key concept box
\newmdenv[
  linecolor=gold,
  linewidth=2pt,
  backgroundcolor=gold!8,
  roundcorner=6pt,
  skipabove=8pt,
  skipbelow=8pt,
  innertopmargin=10pt,
  innerbottommargin=10pt,
  innerleftmargin=12pt,
  innerrightmargin=12pt
]{keybox}

% Warning / important box
\newmdenv[
  linecolor=darkred,
  linewidth=1.5pt,
  backgroundcolor=darkred!5,
  roundcorner=6pt,
  skipabove=6pt,
  skipbelow=6pt,
  innertopmargin=8pt,
  innerbottommargin=8pt,
  innerleftmargin=10pt,
  innerrightmargin=10pt
]{warnbox}

% =====================================================
% DOCUMENT START
% =====================================================
\begin{document}

% =====================================================
% TITLE PAGE
% =====================================================
\begin{titlepage}
  \centering
  \vspace*{2cm}

  {\color{darkblue}\Huge\bfseries YARP RV32I}

  \vspace{0.5cm}

  {\color{accentblue}\LARGE Design \& Verification Guide}

  \vspace{2cm}

  \begin{center}
    \colorbox{lightblue}{
      \parbox{0.8\textwidth}{
        \centering
        \vspace{1cm}
        {\color{darkblue}\Large\bfseries RISC-V RV32I Processor}

        \vspace{0.5cm}

        {\color{accentblue}\normalsize From ISA Fundamentals to Hardware Implementation}
        \vspace{1cm}
      }
    }
  \end{center}

  \vspace{3cm}

  {\color{gold}\large\bfseries Key Topics:}
  \begin{itemize}
    \item[\color{purple}\checkmark] Computer Architecture Fundamentals
    \item[\color{purple}\checkmark] RISC-V RV32I Instruction Set Architecture
    \item[\color{purple}\checkmark] Instruction Formats (R, I, S, B, U, J)
    \item[\color{purple}\checkmark] Complete ISA Reference
    \item[\color{purple}\checkmark] Programmer's Model \& Register File
    \item[\color{purple}\checkmark] YARP Processor Microarchitecture
  \end{itemize}

  \vfill

  {\color{darkgreen}\normalsize Last Updated: \today}

\end{titlepage}

% =====================================================
% TABLE OF CONTENTS
% =====================================================
\tableofcontents
\newpage

% =====================================================
% CHAPTER 1: COMPUTER ARCHITECTURE
% =====================================================
\chapter{Computer Architecture}

\section{What Is Computer Architecture?}

\textbf{Computer architecture} describes the structure and behaviour of a computer system at a level visible to a programmer or hardware designer. It defines \textit{how} a processor executes instructions, \textit{how} it stores and retrieves data, and \textit{how} the various internal components interact.

At the highest level, every computer --- regardless of its complexity --- follows a simple cycle:

\begin{enumerate}
  \item \textbf{\color{purple}Fetch} --- retrieve the next instruction from memory.
  \item \textbf{\color{purple}Decode} --- interpret the instruction to determine what operation is required.
  \item \textbf{\color{purple}Execute} --- carry out the operation (arithmetic, load/store, branch, etc.).
  \item \textbf{\color{purple}Write-back} --- store the result in a register or memory.
\end{enumerate}

This is commonly known as the \textbf{Fetch--Decode--Execute} cycle.

\begin{keybox}
\textbf{\color{gold}Key Concept:}
A processor's \textbf{Instruction Set Architecture (ISA)} is the contract between software and hardware. It specifies the instructions the processor can execute, the registers available to the programmer, and the memory model. The ISA is architecture; the circuits that implement it are \textbf{microarchitecture}.
\end{keybox}

% ─── Von Neumann ──────────────────────────────────────
\section{Von Neumann Architecture}

The \textbf{Von Neumann architecture} (proposed by John von Neumann in 1945) is the foundational model for nearly all modern computers. Its key characteristic is a \textit{single, shared memory} that stores both \textbf{instructions} and \textbf{data}.

\begin{center}
  \includegraphics[width=0.85\textwidth]{images/von_neumann.png}
\end{center}

\subsection{Key Properties}

\begin{itemize}
  \item \textbf{\color{purple}Single memory space:} program code and data occupy the same address range. The CPU fetches instructions and reads/writes data through the \textbf{same bus}.
  \item \textbf{\color{purple}Sequential execution:} one instruction is fetched, decoded, and executed at a time (in the simplest implementation).
  \item \textbf{\color{purple}Stored-program concept:} the program itself resides in modifiable memory, making general-purpose computing possible.
\end{itemize}

\begin{warnbox}
\textbf{\color{darkred}Von Neumann Bottleneck:}
Because instructions and data share a single bus, the CPU can either fetch an instruction \textit{or} access data in a given cycle --- never both simultaneously. This creates a bandwidth bottleneck that limits performance.
\end{warnbox}

% ─── Harvard ──────────────────────────────────────────
\section{Harvard Architecture}

The \textbf{Harvard architecture} solves the Von Neumann bottleneck by providing \textbf{separate memory spaces and buses} for instructions and data.

\begin{center}
  \includegraphics[width=0.8\textwidth]{images/harvard.png}
\end{center}

\subsection{Key Properties}

\begin{itemize}
  \item \textbf{\color{purple}Separate memories:} instruction memory and data memory are physically distinct, each with its own address bus and data bus.
  \item \textbf{\color{purple}Simultaneous access:} the CPU can fetch the next instruction while reading/writing data in the same clock cycle --- doubling effective bandwidth.
  \item \textbf{\color{purple}Common in embedded and RISC:} many modern processors (including RISC-V implementations) use a \textbf{Modified Harvard} architecture: separate caches for instruction and data (L1 I-cache / D-cache) backed by a unified main memory.
\end{itemize}

\begin{infobox}
\textbf{\color{accentblue}RISC-V and Harvard:}
Most practical RISC-V cores adopt a \textbf{Modified Harvard} approach --- split instruction and data caches for performance, but a unified physical memory so that programs can be loaded and modified conventionally.
\end{infobox}

% ─── Inside the CPU ─────────────────────────────────
\section{Inside the CPU}

A CPU is built from several cooperating blocks. The diagram below shows the main components relevant to a simple RV32I implementation:

\begin{center}
  \includegraphics[width=0.78\textwidth]{images/cpu_block.png}
\end{center}

\begin{itemize}
  \item \textbf{\color{purple}Program Counter (PC):} holds the address of the \textit{current} instruction. After each fetch, the PC is incremented by 4 (the size of one RV32I instruction) or updated by a branch/jump.
  \item \textbf{\color{purple}Instruction Decoder:} breaks the 32-bit instruction word into fields (opcode, funct3, funct7, register addresses, immediate value) and determines the required operation.
  \item \textbf{\color{purple}Register File:} a bank of 32 general-purpose registers (x0--x31), each 32 bits wide. The file typically has two read ports and one write port.
  \item \textbf{\color{purple}ALU (Arithmetic Logic Unit):} performs arithmetic (add, subtract) and logic (AND, OR, XOR, shift, compare) operations.
  \item \textbf{\color{purple}Immediate Generator:} extracts and sign-extends the immediate value embedded in the instruction, rearranging the bits depending on the instruction format.
  \item \textbf{\color{purple}Control Unit:} generates control signals that steer multiplexers, the ALU operation selector, register write-enable, memory read/write, and branch decision logic.
\end{itemize}

% ─── RISC vs CISC ─────────────────────────────────────
\section{RISC vs.\ CISC}

Processors are broadly classified into two design philosophies:

\begin{center}
\renewcommand{\arraystretch}{1.35}
\begin{tabular}{p{3.5cm} p{5.8cm} p{5.8cm}}
\toprule
\textbf{\color{purple}Aspect} & \textbf{\color{rtypecolor}CISC} \newline \footnotesize (Complex Instruction Set) & \textbf{\color{itypecolor}RISC} \newline \footnotesize (Reduced Instruction Set) \\
\midrule
Instruction count   & Many (hundreds)                & Few (typically $< 50$ base) \\
Instruction length   & Variable (1--15 bytes on x86)  & Fixed (4 bytes in RV32I) \\
Execution time       & Multi-cycle per instruction     & Ideally 1 cycle per instruction \\
Memory access        & Instructions can operate directly on memory & Load/Store architecture: only \texttt{LW}/\texttt{SW} access memory \\
Complexity           & Complex hardware decoder        & Simple decoder, more compiler work \\
Examples             & x86, x86-64                     & RISC-V, ARM, MIPS \\
\bottomrule
\end{tabular}
\end{center}

\begin{keybox}
\textbf{\color{gold}Why RISC-V?}
RISC-V is an \textbf{open-standard} ISA released by UC Berkeley in 2010. It is modular (start with RV32I, add M, A, F, D, C extensions as needed), royalty-free, and designed for clean hardware implementation --- making it an ideal educational and commercial platform.
\end{keybox}

\newpage

% =====================================================
% CHAPTER 2: THE RV32I ISA
% =====================================================
\chapter{The RV32I Instruction Set Architecture}

\section{Introduction to RISC-V}

\textbf{RISC-V} (pronounced ``risk-five'') is a free and open Instruction Set Architecture originally developed at the University of California, Berkeley. Unlike proprietary ISAs (x86, ARM), \textbf{anyone} can implement a RISC-V processor without licensing fees.

The ISA is \textbf{modular}: a small mandatory base (\texttt{RV32I}) provides a complete integer instruction set, and optional standard extensions add functionality:

\begin{center}
\renewcommand{\arraystretch}{1.3}
\begin{tabular}{c l l}
\toprule
\textbf{\color{purple}Extension} & \textbf{\color{purple}Name} & \textbf{\color{purple}Description} \\
\midrule
\texttt{I} & Base Integer        & 32-bit integer arithmetic, loads, stores, branches, jumps \\
\texttt{M} & Multiply / Divide   & Hardware multiply and divide instructions \\
\texttt{A} & Atomic              & Atomic read-modify-write for multi-core \\
\texttt{F} & Single-Float        & IEEE 754 single-precision floating point \\
\texttt{D} & Double-Float        & IEEE 754 double-precision floating point \\
\texttt{C} & Compressed          & 16-bit compressed instruction encoding \\
\texttt{Zicsr} & CSR              & Control and Status Register instructions \\
\bottomrule
\end{tabular}
\end{center}

\vspace{0.3cm}

Throughout this guide we focus exclusively on the \textbf{RV32I} base --- the minimum set that every RISC-V implementation must support.

\section{RV32I Overview}

The RV32I base defines:

\begin{itemize}
  \item \textbf{\color{purple}32 general-purpose registers} (\texttt{x0}--\texttt{x31}), each 32 bits wide. \texttt{x0} is hardwired to zero.
  \item \textbf{\color{purple}A 32-bit Program Counter (PC)}.
  \item \textbf{\color{purple}Fixed 32-bit instruction width} (4 bytes, naturally aligned).
  \item \textbf{\color{purple}A load/store architecture:} only load and store instructions access memory; all other operations work on registers.
  \item \textbf{\color{purple}40 base instructions} organized into 6 encoding formats.
  \item \textbf{\color{purple}Little-endian} byte ordering (with optional big-endian).
  \item \textbf{\color{purple}32-bit flat address space} ($2^{32} = 4\,\text{GiB}$).
\end{itemize}

% ─── Programmer's Model ──────────────────────────────
\section{Programmer's Model}

\subsection{The Register File}

The RV32I register file contains \textbf{32 integer registers}, each 32 bits wide, designated \texttt{x0} through \texttt{x31}. In addition, each register has a \textit{conventional ABI (Application Binary Interface) name} used by compilers and assembly programmers.

\begin{center}
  \includegraphics[width=0.55\textwidth]{images/register_file.png}
\end{center}

\begin{keybox}
\textbf{\color{gold}Special Register --- \texttt{x0} (zero):}
Register \texttt{x0} is \textbf{hardwired to the value 0}. Any write to \texttt{x0} is silently discarded. This simplifies many operations --- for example, \texttt{addi x5, x0, 42} loads the constant 42 into \texttt{x5} without needing a separate ``load immediate'' instruction.
\end{keybox}

\subsection{Register Convention (ABI)}

\begin{center}
\renewcommand{\arraystretch}{1.2}
\footnotesize
\begin{tabular}{c c l l}
\toprule
\textbf{\color{purple}Register} & \textbf{\color{purple}ABI Name} & \textbf{\color{purple}Role} & \textbf{\color{purple}Saver} \\
\midrule
\texttt{x0}       & \texttt{zero} & Hard-wired zero               & --- \\
\texttt{x1}       & \texttt{ra}   & Return address                & Caller \\
\texttt{x2}       & \texttt{sp}   & Stack pointer                 & Callee \\
\texttt{x3}       & \texttt{gp}   & Global pointer                & --- \\
\texttt{x4}       & \texttt{tp}   & Thread pointer                & --- \\
\texttt{x5--x7}   & \texttt{t0--t2} & Temporaries                 & Caller \\
\texttt{x8}       & \texttt{s0/fp}  & Saved / frame pointer       & Callee \\
\texttt{x9}       & \texttt{s1}     & Saved register              & Callee \\
\texttt{x10--x11} & \texttt{a0--a1} & Function args / return vals & Caller \\
\texttt{x12--x17} & \texttt{a2--a7} & Function arguments          & Caller \\
\texttt{x18--x27} & \texttt{s2--s11}& Saved registers             & Callee \\
\texttt{x28--x31} & \texttt{t3--t6} & Temporaries                 & Caller \\
\bottomrule
\end{tabular}
\end{center}

\begin{infobox}
\textbf{\color{accentblue}Caller-saved vs.\ Callee-saved:}
\textbf{Caller-saved} registers may be overwritten by any function call; the caller must save them if their values are needed later. \textbf{Callee-saved} registers must be preserved across function calls; if a function modifies them, it must save and restore them.
\end{infobox}

\subsection{Program Counter}

The \textbf{Program Counter (PC)} is a dedicated 32-bit register that holds the address of the \textit{currently executing instruction}. After each instruction, the PC is updated:
\begin{itemize}
  \item \textbf{Sequential:} $\text{PC} \leftarrow \text{PC} + 4$ (next instruction).
  \item \textbf{Branch taken:} $\text{PC} \leftarrow \text{PC} + \text{offset}$ (relative jump).
  \item \textbf{Jump:} $\text{PC} \leftarrow \text{target address}$ (absolute or register-based).
\end{itemize}

\subsection{Memory Model}

RV32I uses a \textbf{flat, byte-addressable} 32-bit address space. Memory is accessed only via load and store instructions.

\begin{center}
  \includegraphics[width=0.45\textwidth]{images/memory_layout.png}
\end{center}

\newpage

% ─── Instruction Formats ─────────────────────────────
\section{Instruction Formats}

Every RV32I instruction is exactly \textbf{32 bits} wide. The ISA defines \textbf{six encoding formats} that specify how the 32 bits are partitioned into fields:

\begin{center}
  \includegraphics[width=0.95\textwidth]{images/instruction_formats.png}
\end{center}

All formats share a common 7-bit \texttt{opcode} field in bits [6:0], which determines the general category of the instruction. The remaining 25 bits carry register addresses, function selectors, and/or immediate values.

\begin{center}
  \includegraphics[width=0.82\textwidth]{images/type_overview.png}
\end{center}

\newpage

% ─── R-Type ───────────────────────────────────────────
\section{R-Type Instructions (Register--Register)}

\begin{infobox}
\textbf{\color{rtypecolor}\Large R-Type} \hfill \texttt{opcode = 0110011}\\[6pt]
R-type instructions perform operations between \textbf{two source registers} (\texttt{rs1}, \texttt{rs2}) and write the result to a \textbf{destination register} (\texttt{rd}). No immediate value is used.
\end{infobox}

\subsection{Format}

\begin{center}
\renewcommand{\arraystretch}{1.4}
\begin{tabular}{|>{\centering\arraybackslash\columncolor{funct7bg}}p{2.2cm}|>{\centering\arraybackslash\columncolor{rs2bg}}p{1.6cm}|>{\centering\arraybackslash\columncolor{rs1bg}}p{1.6cm}|>{\centering\arraybackslash\columncolor{funct3bg}}p{1.5cm}|>{\centering\arraybackslash\columncolor{rdbg}}p{1.6cm}|>{\centering\arraybackslash\columncolor{opcodebg}}p{2.2cm}|}
\hline
\textbf{funct7} & \textbf{rs2} & \textbf{rs1} & \textbf{funct3} & \textbf{rd} & \textbf{opcode} \\
\hline
\footnotesize bits [31:25] & \footnotesize bits [24:20] & \footnotesize bits [19:15] & \footnotesize bits [14:12] & \footnotesize bits [11:7] & \footnotesize bits [6:0] \\
\hline
\footnotesize 7 bits & \footnotesize 5 bits & \footnotesize 5 bits & \footnotesize 3 bits & \footnotesize 5 bits & \footnotesize 7 bits \\
\hline
\end{tabular}
\end{center}

\subsection{Fields}

\begin{itemize}
  \item \textbf{\color{purple}opcode} (7 bits): always \texttt{0110011} for R-type integer operations.
  \item \textbf{\color{purple}rd} (5 bits): destination register that receives the result.
  \item \textbf{\color{purple}funct3} (3 bits): selects the specific operation (ADD/SUB, SLL, SLT, etc.).
  \item \textbf{\color{purple}rs1} (5 bits): first source register.
  \item \textbf{\color{purple}rs2} (5 bits): second source register.
  \item \textbf{\color{purple}funct7} (7 bits): further disambiguates the operation (e.g., \texttt{0000000} for ADD vs.\ \texttt{0100000} for SUB).
\end{itemize}

\subsection{R-Type Instruction Listing}

\begin{center}
\renewcommand{\arraystretch}{1.25}
\begin{tabular}{l l l p{7.5cm}}
\toprule
\textbf{\color{purple}Instruction} & \textbf{\color{purple}funct7} & \textbf{\color{purple}funct3} & \textbf{\color{purple}Operation} \\
\midrule
\texttt{ADD  rd, rs1, rs2}  & \texttt{0000000} & \texttt{000} & $\texttt{rd} \leftarrow \texttt{rs1} + \texttt{rs2}$ \\
\texttt{SUB  rd, rs1, rs2}  & \texttt{0100000} & \texttt{000} & $\texttt{rd} \leftarrow \texttt{rs1} - \texttt{rs2}$ \\
\texttt{SLL  rd, rs1, rs2}  & \texttt{0000000} & \texttt{001} & $\texttt{rd} \leftarrow \texttt{rs1} \ll \texttt{rs2}[4:0]$ (logical left shift) \\
\texttt{SLT  rd, rs1, rs2}  & \texttt{0000000} & \texttt{010} & $\texttt{rd} \leftarrow (\texttt{rs1} <_s \texttt{rs2})\;?\;1:0$ (signed compare) \\
\texttt{SLTU rd, rs1, rs2}  & \texttt{0000000} & \texttt{011} & $\texttt{rd} \leftarrow (\texttt{rs1} <_u \texttt{rs2})\;?\;1:0$ (unsigned compare) \\
\texttt{XOR  rd, rs1, rs2}  & \texttt{0000000} & \texttt{100} & $\texttt{rd} \leftarrow \texttt{rs1} \oplus \texttt{rs2}$ (bitwise XOR) \\
\texttt{SRL  rd, rs1, rs2}  & \texttt{0000000} & \texttt{101} & $\texttt{rd} \leftarrow \texttt{rs1} \gg_\text{logic} \texttt{rs2}[4:0]$ (logical right shift) \\
\texttt{SRA  rd, rs1, rs2}  & \texttt{0100000} & \texttt{101} & $\texttt{rd} \leftarrow \texttt{rs1} \gg_\text{arith} \texttt{rs2}[4:0]$ (arithmetic right shift) \\
\texttt{OR   rd, rs1, rs2}  & \texttt{0000000} & \texttt{110} & $\texttt{rd} \leftarrow \texttt{rs1} \,|\, \texttt{rs2}$ (bitwise OR) \\
\texttt{AND  rd, rs1, rs2}  & \texttt{0000000} & \texttt{111} & $\texttt{rd} \leftarrow \texttt{rs1} \,\&\, \texttt{rs2}$ (bitwise AND) \\
\bottomrule
\end{tabular}
\end{center}

\begin{infobox}
\textbf{\color{accentblue}Note:} ADD and SUB share the same \texttt{funct3 = 000}. They are distinguished by \texttt{funct7}: \texttt{0000000} for ADD, \texttt{0100000} for SUB. The same pattern applies to SRL/SRA (\texttt{funct3 = 101}).
\end{infobox}

\begin{keybox}
\textbf{\color{accentblue}Shift reminder:} A bitwise operation works directly on the 0/1 pattern stored in a register. A shift operation moves that bit pattern left or right by some number of positions.

\begin{itemize}
  \item \textbf{\texttt{SLL} (shift left logical):} moves bits to the left, drops any bits that fall off the left side, and fills the new right-side bits with 0s. It is often used to multiply by $2^n$ when overflow is not a concern.
  \item \textbf{\texttt{SRL} (shift right logical):} moves bits to the right and fills the new left-side bits with 0s. This is useful for unsigned values and for extracting fields from an instruction or address.
  \item \textbf{\texttt{SRA} (shift right arithmetic):} moves bits to the right but copies the old leftmost bit (the sign bit) into the new left-side positions. This preserves the sign of a signed value.
\end{itemize}

There is no separate arithmetic left shift in RV32I: left shifts already insert 0s on the right, so \texttt{SLL} is the left-shift instruction you use.

\textbf{Small example (shown in 8 bits):}
\[
\texttt{00101101} \ll 1 = \texttt{01011010}
\]
\[
\texttt{11110100} \gg_\text{logic} 1 = \texttt{01111010}, \qquad
\texttt{11110100} \gg_\text{arith} 1 = \texttt{11111010}
\]
If \texttt{11110100} represents $-12$ in two's complement, arithmetic right shift keeps the value negative, while logical right shift inserts a 0 on the left and changes the meaning.

\textbf{Typical use cases:} use \texttt{SLL} to scale an index by 2, 4, or 8, use \texttt{SRL} to unpack bitfields, and use \texttt{SRA} when shifting signed values right without losing the sign.
\end{keybox}

\newpage

% ─── I-Type ───────────────────────────────────────────
\section{I-Type Instructions (Immediate)}

\begin{infobox}
\textbf{\color{itypecolor}\Large I-Type} \hfill \texttt{opcodes: 0010011, 0000011, 1100111, 1110011}\\[6pt]
I-type instructions use a \textbf{12-bit signed immediate} value as one operand alongside a source register. This format is used for \textbf{arithmetic with immediates}, \textbf{loads from memory}, \textbf{JALR}, and \textbf{system calls}.
\end{infobox}

\subsection{Format}

\begin{center}
\renewcommand{\arraystretch}{1.4}
\begin{tabular}{|>{\centering\arraybackslash\columncolor{immbg}}p{3.8cm}|>{\centering\arraybackslash\columncolor{rs1bg}}p{1.6cm}|>{\centering\arraybackslash\columncolor{funct3bg}}p{1.5cm}|>{\centering\arraybackslash\columncolor{rdbg}}p{1.6cm}|>{\centering\arraybackslash\columncolor{opcodebg}}p{2.2cm}|}
\hline
\textbf{imm[11:0]} & \textbf{rs1} & \textbf{funct3} & \textbf{rd} & \textbf{opcode} \\
\hline
\footnotesize bits [31:20] & \footnotesize bits [19:15] & \footnotesize bits [14:12] & \footnotesize bits [11:7] & \footnotesize bits [6:0] \\
\hline
\footnotesize 12 bits & \footnotesize 5 bits & \footnotesize 3 bits & \footnotesize 5 bits & \footnotesize 7 bits \\
\hline
\end{tabular}
\end{center}

\subsection{Fields}

\begin{itemize}
  \item \textbf{\color{purple}imm[11:0]} (12 bits): a \textbf{sign-extended} immediate value (range: $-2048$ to $+2047$).
  \item \textbf{\color{purple}rs1} (5 bits): source register.
  \item \textbf{\color{purple}funct3} (3 bits): selects the operation.
  \item \textbf{\color{purple}rd} (5 bits): destination register.
  \item \textbf{\color{purple}opcode} (7 bits): identifies the instruction sub-category.
\end{itemize}

\subsection{I-Type Arithmetic Instructions (\texttt{opcode = 0010011})}

\begin{center}
\renewcommand{\arraystretch}{1.25}
\begin{tabular}{l l p{8cm}}
\toprule
\textbf{\color{purple}Instruction} & \textbf{\color{purple}funct3} & \textbf{\color{purple}Operation} \\
\midrule
\texttt{ADDI  rd, rs1, imm}  & \texttt{000} & $\texttt{rd} \leftarrow \texttt{rs1} + \text{sext}(\text{imm})$ \\
\texttt{SLTI  rd, rs1, imm}  & \texttt{010} & $\texttt{rd} \leftarrow (\texttt{rs1} <_s \text{sext}(\text{imm}))\;?\;1:0$ \\
\texttt{SLTIU rd, rs1, imm}  & \texttt{011} & $\texttt{rd} \leftarrow (\texttt{rs1} <_u \text{sext}(\text{imm}))\;?\;1:0$ \\
\texttt{XORI  rd, rs1, imm}  & \texttt{100} & $\texttt{rd} \leftarrow \texttt{rs1} \oplus \text{sext}(\text{imm})$ \\
\texttt{ORI   rd, rs1, imm}  & \texttt{110} & $\texttt{rd} \leftarrow \texttt{rs1} \,|\, \text{sext}(\text{imm})$ \\
\texttt{ANDI  rd, rs1, imm}  & \texttt{111} & $\texttt{rd} \leftarrow \texttt{rs1} \,\&\, \text{sext}(\text{imm})$ \\
\texttt{SLLI  rd, rs1, shamt}& \texttt{001} & $\texttt{rd} \leftarrow \texttt{rs1} \ll \text{shamt}$ \hfill {\scriptsize imm[11:5]=\texttt{0000000}} \\
\texttt{SRLI  rd, rs1, shamt}& \texttt{101} & $\texttt{rd} \leftarrow \texttt{rs1} \gg_\text{logic} \text{shamt}$ \hfill {\scriptsize imm[11:5]=\texttt{0000000}} \\
\texttt{SRAI  rd, rs1, shamt}& \texttt{101} & $\texttt{rd} \leftarrow \texttt{rs1} \gg_\text{arith} \text{shamt}$ \hfill {\scriptsize imm[11:5]=\texttt{0100000}} \\
\bottomrule
\end{tabular}
\end{center}

\begin{infobox}
\textbf{\color{accentblue}Shift immediates:}
For SLLI, SRLI, and SRAI, only the lower 5 bits of the immediate (\texttt{shamt} = shift amount, bits [4:0]) specify the shift distance (0--31). The upper bits (imm[11:5]) act like \texttt{funct7} to distinguish SRLI from SRAI.
\end{infobox}

\begin{keybox}
\textbf{\color{gold}How to think about I-Type arithmetic:}
Most I-type arithmetic instructions mean ``take one register value and combine it with a small constant.'' This is why they are heavily used for loop counters, stack-pointer updates, small address adjustments, and bit masks.

\begin{itemize}
  \item \textbf{\texttt{ADDI}} is the workhorse for incrementing and decrementing values.
  \item There is no separate \texttt{SUBI} instruction in RV32I: subtracting a constant is just \texttt{ADDI} with a negative immediate.
  \item \textbf{\texttt{ANDI}/\texttt{ORI}/\texttt{XORI}} perform bitwise operations between \texttt{rs1} and the immediate, which makes them useful for masking or setting selected bits.
\end{itemize}

\textbf{Small example:} \texttt{addi x5, x5, -4} subtracts 4 from \texttt{x5}.
\end{keybox}

\subsection{I-Type Load Instructions (\texttt{opcode = 0000011})}

Load instructions read data from memory at address $\texttt{rs1} + \text{sext}(\text{imm})$ and place it in \texttt{rd}.

\begin{center}
\renewcommand{\arraystretch}{1.25}
\begin{tabular}{l l p{8cm}}
\toprule
\textbf{\color{purple}Instruction} & \textbf{\color{purple}funct3} & \textbf{\color{purple}Operation} \\
\midrule
\texttt{LB  rd, imm(rs1)}  & \texttt{000} & Load byte (8-bit), \textbf{sign-extend} to 32 bits $\rightarrow$ \texttt{rd} \\
\texttt{LH  rd, imm(rs1)}  & \texttt{001} & Load halfword (16-bit), \textbf{sign-extend} to 32 bits $\rightarrow$ \texttt{rd} \\
\texttt{LW  rd, imm(rs1)}  & \texttt{010} & Load word (32-bit) $\rightarrow$ \texttt{rd} \\
\texttt{LBU rd, imm(rs1)}  & \texttt{100} & Load byte (8-bit), \textbf{zero-extend} to 32 bits $\rightarrow$ \texttt{rd} \\
\texttt{LHU rd, imm(rs1)}  & \texttt{101} & Load halfword (16-bit), \textbf{zero-extend} to 32 bits $\rightarrow$ \texttt{rd} \\
\bottomrule
\end{tabular}
\end{center}

\begin{infobox}
\textbf{\color{accentblue}Sign-extension vs.\ Zero-extension:}
\texttt{LB}/\texttt{LH} are \textbf{signed}: the MSB of the loaded value is replicated to fill the upper bits, preserving the sign of the original value. \texttt{LBU}/\texttt{LHU} are \textbf{unsigned}: the upper bits are filled with zeros. This is critical when loading narrow data into a 32-bit register.
\end{infobox}

\begin{keybox}
\textbf{\color{gold}Width reminder:}
In RV32I, a \textbf{byte} is 8 bits, a \textbf{halfword} is 16 bits, and a \textbf{word} is 32 bits. The load instruction tells the processor both \textbf{how many bytes to read} and \textbf{how to fill the rest of the destination register}.

\begin{itemize}
  \item Use \texttt{LB}/\texttt{LBU} for 8-bit values such as characters or packed flags.
  \item Use \texttt{LH}/\texttt{LHU} for 16-bit fields.
  \item Use \texttt{LW} when the full 32-bit word is needed.
\end{itemize}

In a simple RV32I implementation, it is safest to think of halfwords as 2-byte aligned and words as 4-byte aligned.
\end{keybox}

\newpage

% ─── S-Type ───────────────────────────────────────────
\section{S-Type Instructions (Store)}

\begin{infobox}
\textbf{\color{stypecolor}\Large S-Type} \hfill \texttt{opcode = 0100011}\\[6pt]
S-type instructions \textbf{store} data from a register (\texttt{rs2}) into memory. The memory address is computed as $\texttt{rs1} + \text{sext}(\text{imm})$. There is \textbf{no destination register} --- the immediate is split across two fields.
\end{infobox}

\subsection{Format}

\begin{center}
\renewcommand{\arraystretch}{1.4}
\begin{tabular}{|>{\centering\arraybackslash\columncolor{immbg}}p{2.2cm}|>{\centering\arraybackslash\columncolor{rs2bg}}p{1.6cm}|>{\centering\arraybackslash\columncolor{rs1bg}}p{1.6cm}|>{\centering\arraybackslash\columncolor{funct3bg}}p{1.5cm}|>{\centering\arraybackslash\columncolor{immbg}}p{1.6cm}|>{\centering\arraybackslash\columncolor{opcodebg}}p{2.2cm}|}
\hline
\textbf{imm[11:5]} & \textbf{rs2} & \textbf{rs1} & \textbf{funct3} & \textbf{imm[4:0]} & \textbf{opcode} \\
\hline
\footnotesize bits [31:25] & \footnotesize bits [24:20] & \footnotesize bits [19:15] & \footnotesize bits [14:12] & \footnotesize bits [11:7] & \footnotesize bits [6:0] \\
\hline
\footnotesize 7 bits & \footnotesize 5 bits & \footnotesize 5 bits & \footnotesize 3 bits & \footnotesize 5 bits & \footnotesize 7 bits \\
\hline
\end{tabular}
\end{center}

\subsection{Why Is the Immediate Split?}

The S-type immediate occupies the same bit positions as R-type's \texttt{funct7} (bits [31:25]) and \texttt{rd} (bits [11:7]). This keeps the \texttt{rs1} and \texttt{rs2} fields in exactly the same position across R, I, and S formats, simplifying the hardware register-file read logic.

\subsection{S-Type Instruction Listing}

\begin{center}
\renewcommand{\arraystretch}{1.25}
\begin{tabular}{l l p{8cm}}
\toprule
\textbf{\color{purple}Instruction} & \textbf{\color{purple}funct3} & \textbf{\color{purple}Operation} \\
\midrule
\texttt{SB rs2, imm(rs1)} & \texttt{000} & Store byte: $\text{Mem}[\texttt{rs1}+\text{sext}(\text{imm})] \leftarrow \texttt{rs2}[7:0]$ \\
\texttt{SH rs2, imm(rs1)} & \texttt{001} & Store halfword: $\text{Mem}[\texttt{rs1}+\text{sext}(\text{imm})] \leftarrow \texttt{rs2}[15:0]$ \\
\texttt{SW rs2, imm(rs1)} & \texttt{010} & Store word: $\text{Mem}[\texttt{rs1}+\text{sext}(\text{imm})] \leftarrow \texttt{rs2}[31:0]$ \\
\bottomrule
\end{tabular}
\end{center}

\begin{keybox}
\textbf{\color{gold}Store reminder:}
Stores move data \textbf{from a register into memory}, so there is nothing to sign-extend or zero-extend at store time.

\begin{itemize}
  \item \texttt{SB} writes only the lower 8 bits of \texttt{rs2}.
  \item \texttt{SH} writes only the lower 16 bits of \texttt{rs2}.
  \item \texttt{SW} writes the full 32 bits of \texttt{rs2}.
\end{itemize}

That is why RV32I has \texttt{LBU}/\texttt{LHU} for unsigned loads but no ``unsigned store'' instructions.

\textbf{Small example:} if \texttt{x6 = 0x12345678}, then \texttt{sb x6, 0(x5)} stores only \texttt{0x78} to memory.
\end{keybox}

\newpage

% ─── B-Type ───────────────────────────────────────────
\section{B-Type Instructions (Branch)}

\begin{infobox}
\textbf{\color{btypecolor}\Large B-Type} \hfill \texttt{opcode = 1100011}\\[6pt]
B-type instructions implement \textbf{conditional branches}. They compare two registers (\texttt{rs1}, \texttt{rs2}) and, if the condition is true, jump to $\text{PC} + \text{sext}(\text{offset})$. Otherwise, execution continues at $\text{PC} + 4$. The offset is a \textbf{13-bit signed} value with the least significant bit always 0 (instructions are 4-byte aligned, offsets are multiples of 2).
\end{infobox}

\subsection{Format}

\begin{center}
\renewcommand{\arraystretch}{1.4}
\footnotesize
\begin{tabular}{|>{\centering\arraybackslash\columncolor{immbg}}p{0.8cm}|>{\centering\arraybackslash\columncolor{immbg}}p{1.8cm}|>{\centering\arraybackslash\columncolor{rs2bg}}p{1.4cm}|>{\centering\arraybackslash\columncolor{rs1bg}}p{1.4cm}|>{\centering\arraybackslash\columncolor{funct3bg}}p{1.2cm}|>{\centering\arraybackslash\columncolor{immbg}}p{1.2cm}|>{\centering\arraybackslash\columncolor{immbg}}p{0.8cm}|>{\centering\arraybackslash\columncolor{opcodebg}}p{1.8cm}|}
\hline
\textbf{[12]} & \textbf{[10:5]} & \textbf{rs2} & \textbf{rs1} & \textbf{f3} & \textbf{[4:1]} & \textbf{[11]} & \textbf{opcode} \\
\hline
\tiny bit 31 & \tiny bits [30:25] & \tiny bits [24:20] & \tiny bits [19:15] & \tiny bits [14:12] & \tiny bits [11:8] & \tiny bit 7 & \tiny bits [6:0] \\
\hline
\tiny 1 & \tiny 6 & \tiny 5 & \tiny 5 & \tiny 3 & \tiny 4 & \tiny 1 & \tiny 7 \\
\hline
\end{tabular}
\end{center}

\subsection{Immediate Reconstruction}

The 13-bit branch offset is assembled from the scattered fields:
$$\text{offset} = \{\text{imm}[12], \text{imm}[11], \text{imm}[10:5], \text{imm}[4:1], 0\}$$
Bit [0] is always 0, giving a range of $\pm 4\,\text{KiB}$ ($-4096$ to $+4094$).

\subsection{B-Type Instruction Listing}

\begin{center}
\renewcommand{\arraystretch}{1.25}
\begin{tabular}{l l p{8cm}}
\toprule
\textbf{\color{purple}Instruction} & \textbf{\color{purple}funct3} & \textbf{\color{purple}Condition (branch if true)} \\
\midrule
\texttt{BEQ  rs1, rs2, offset} & \texttt{000} & $\texttt{rs1} = \texttt{rs2}$ \hfill (equal) \\
\texttt{BNE  rs1, rs2, offset} & \texttt{001} & $\texttt{rs1} \neq \texttt{rs2}$ \hfill (not equal) \\
\texttt{BLT  rs1, rs2, offset} & \texttt{100} & $\texttt{rs1} <_s \texttt{rs2}$ \hfill (signed less than) \\
\texttt{BGE  rs1, rs2, offset} & \texttt{101} & $\texttt{rs1} \geq_s \texttt{rs2}$ \hfill (signed greater or equal) \\
\texttt{BLTU rs1, rs2, offset} & \texttt{110} & $\texttt{rs1} <_u \texttt{rs2}$ \hfill (unsigned less than) \\
\texttt{BGEU rs1, rs2, offset} & \texttt{111} & $\texttt{rs1} \geq_u \texttt{rs2}$ \hfill (unsigned greater or equal) \\
\bottomrule
\end{tabular}
\end{center}

\begin{infobox}
\textbf{\color{accentblue}No BLE / BGT?}
RISC-V omits ``branch-if-less-or-equal'' and ``branch-if-greater-than'' because they can be synthesized by swapping operands:
\begin{itemize}
  \item \texttt{BLE rs1, rs2, offset} $\equiv$ \texttt{BGE rs2, rs1, offset}
  \item \texttt{BGT rs1, rs2, offset} $\equiv$ \texttt{BLT rs2, rs1, offset}
\end{itemize}
This keeps the hardware simple without losing expressiveness.
\end{infobox}

\begin{keybox}
\textbf{\color{gold}Branch reminder:}
A branch does not contain a full absolute destination address. Instead, it says ``if this comparison is true, add this signed offset to the current \texttt{PC}.'' If the comparison is false, execution simply continues at the next instruction (\texttt{PC + 4}).

Because bit 0 of the encoded offset is always 0, the target moves in steps of 2 bytes rather than 1 byte.

\textbf{Small example:} if the current \texttt{PC} is \texttt{0x1000} and the branch offset is \texttt{+16}, the branch target is \texttt{0x1010}. If the condition is false, execution continues at \texttt{0x1004}.
\end{keybox}

\newpage

% ─── U-Type ───────────────────────────────────────────
\section{U-Type Instructions (Upper Immediate)}

\begin{infobox}
\textbf{\color{utypecolor}\Large U-Type} \hfill \texttt{opcodes: 0110111 (LUI), 0010111 (AUIPC)}\\[6pt]
U-type instructions carry a \textbf{20-bit upper immediate} that is placed into bits [31:12] of the destination register. The lower 12 bits are set to zero. This format is used to construct large constants or PC-relative addresses.
\end{infobox}

\subsection{Format}

\begin{center}
\renewcommand{\arraystretch}{1.4}
\begin{tabular}{|>{\centering\arraybackslash\columncolor{immbg}}p{6.5cm}|>{\centering\arraybackslash\columncolor{rdbg}}p{1.6cm}|>{\centering\arraybackslash\columncolor{opcodebg}}p{2.2cm}|}
\hline
\textbf{imm[31:12]} & \textbf{rd} & \textbf{opcode} \\
\hline
\footnotesize bits [31:12] & \footnotesize bits [11:7] & \footnotesize bits [6:0] \\
\hline
\footnotesize 20 bits & \footnotesize 5 bits & \footnotesize 7 bits \\
\hline
\end{tabular}
\end{center}

\subsection{U-Type Instruction Listing}

\begin{center}
\renewcommand{\arraystretch}{1.25}
\begin{tabular}{l l p{8.5cm}}
\toprule
\textbf{\color{purple}Instruction} & \textbf{\color{purple}Opcode} & \textbf{\color{purple}Operation} \\
\midrule
\texttt{LUI   rd, imm} & \texttt{0110111} & $\texttt{rd} \leftarrow \text{imm}[31:12] \,\|\, 12'b0$ \newline Load Upper Immediate: places the 20-bit immediate in the upper 20 bits of \texttt{rd}, zeroing the lower 12 bits. \\[6pt]
\texttt{AUIPC rd, imm} & \texttt{0010111} & $\texttt{rd} \leftarrow \text{PC} + (\text{imm}[31:12] \,\|\, 12'b0)$ \newline Add Upper Immediate to PC: adds the shifted immediate to the current PC and stores the result in \texttt{rd}. Used for PC-relative addressing of data/code. \\
\bottomrule
\end{tabular}
\end{center}

\begin{keybox}
\textbf{\color{gold}Building 32-bit Constants:}
Since immediates in I-type are only 12 bits, you need \textbf{two instructions} to load a full 32-bit constant:
\begin{enumerate}
  \item \texttt{LUI rd, upper\_20} --- sets bits [31:12]
  \item \texttt{ADDI rd, rd, lower\_12} --- sets bits [11:0]
\end{enumerate}
Example: load \texttt{0xDEADBEEF} into \texttt{x5}:
\begin{lstlisting}
  lui  x5, 0xDEADC     # x5 = 0xDEADC000
  addi x5, x5, 0xEEF   # x5 = 0xDEADBEEF  (note: sign-extension adjustment)
\end{lstlisting}
\vspace{-0.3cm}
\end{keybox}

\begin{infobox}
\textbf{\color{accentblue}LUI vs.\ AUIPC mental model:}
\texttt{LUI} builds a value from the immediate itself, while \texttt{AUIPC} builds a value \textbf{relative to the current \texttt{PC}}.

Use \texttt{LUI} when you want a large constant in a register. Use \texttt{AUIPC} when you want a base address for nearby code or data.

That is why \texttt{AUIPC} often appears with \texttt{JALR}, loads, or stores in PC-relative code sequences.
\end{infobox}

\newpage

% ─── J-Type ───────────────────────────────────────────
\section{J-Type Instructions (Jump)}

\begin{infobox}
\textbf{\color{jtypecolor}\Large J-Type} \hfill \texttt{opcode = 1101111 (JAL)}\\[6pt]
J-type provides an unconditional jump with a \textbf{21-bit signed offset} (range: $\pm 1\,\text{MiB}$). The return address ($\text{PC}+4$) is saved in \texttt{rd}, enabling function calls.
\end{infobox}

\subsection{Format}

\begin{center}
\renewcommand{\arraystretch}{1.4}
\footnotesize
\begin{tabular}{|>{\centering\arraybackslash\columncolor{immbg}}p{0.8cm}|>{\centering\arraybackslash\columncolor{immbg}}p{3.0cm}|>{\centering\arraybackslash\columncolor{immbg}}p{0.8cm}|>{\centering\arraybackslash\columncolor{immbg}}p{2.2cm}|>{\centering\arraybackslash\columncolor{rdbg}}p{1.4cm}|>{\centering\arraybackslash\columncolor{opcodebg}}p{1.8cm}|}
\hline
\textbf{[20]} & \textbf{[10:1]} & \textbf{[11]} & \textbf{[19:12]} & \textbf{rd} & \textbf{opcode} \\
\hline
\tiny bit 31 & \tiny bits [30:21] & \tiny bit 20 & \tiny bits [19:12] & \tiny bits [11:7] & \tiny bits [6:0] \\
\hline
\tiny 1 & \tiny 10 & \tiny 1 & \tiny 8 & \tiny 5 & \tiny 7 \\
\hline
\end{tabular}
\end{center}

\subsection{Immediate Reconstruction}

The 21-bit offset is assembled as:
$$\text{offset} = \{\text{imm}[20], \text{imm}[19:12], \text{imm}[11], \text{imm}[10:1], 0\}$$

\subsection{Jump Instructions}

\begin{center}
\renewcommand{\arraystretch}{1.25}
\begin{tabular}{l l p{8cm}}
\toprule
\textbf{\color{purple}Instruction} & \textbf{\color{purple}Type} & \textbf{\color{purple}Operation} \\
\midrule
\texttt{JAL  rd, offset}       & J-type & $\texttt{rd} \leftarrow \text{PC}+4$; \quad $\text{PC} \leftarrow \text{PC} + \text{sext}(\text{offset})$ \newline Jump And Link: saves return address in \texttt{rd}, then jumps to PC-relative target. \\[6pt]
\texttt{JALR rd, rs1, offset}  & I-type & $\texttt{rd} \leftarrow \text{PC}+4$; \quad $\text{PC} \leftarrow (\texttt{rs1} + \text{sext}(\text{imm})) \;\&\; \sim1$ \newline Jump And Link Register: jumps to an address computed from a register, clearing bit 0. Used for indirect jumps and function returns. \\
\bottomrule
\end{tabular}
\end{center}

\begin{infobox}
\textbf{\color{accentblue}Function call / return pattern:}
\begin{itemize}
  \item \textbf{Call:} \texttt{jal ra, function\_label} --- jumps to function, saves return address in \texttt{ra} (\texttt{x1}).
  \item \textbf{Return:} \texttt{jalr zero, ra, 0} --- jumps to address in \texttt{ra}, discards link (writes to \texttt{x0}).
\end{itemize}
\end{infobox}

\begin{keybox}
\textbf{\color{gold}What does ``link'' mean?}
The \textbf{link} is just the return address, usually \texttt{PC + 4}. Saving it in \texttt{ra} lets the called function know where to come back after it finishes.

If you do not need a return address, you can discard the link by writing it to \texttt{x0}. That turns \texttt{JAL} into a plain unconditional jump.

\begin{itemize}
  \item \texttt{jal ra, function\_label} --- function call
  \item \texttt{jal x0, label} --- unconditional jump without saving a return address
\end{itemize}

\texttt{JALR} is the register-based version, which is why returns are commonly written as \texttt{jalr x0, ra, 0}.
\end{keybox}

\newpage

% ─── System Instructions ─────────────────────────────
\section{System Instructions}

\begin{infobox}
\textbf{\color{gold}\Large System} \hfill \texttt{opcode = 1110011}\\[6pt]
The RV32I base includes two system instructions used to request services from the execution environment (OS or debugger).
\end{infobox}

\begin{center}
\renewcommand{\arraystretch}{1.25}
\begin{tabular}{l p{10cm}}
\toprule
\textbf{\color{purple}Instruction} & \textbf{\color{purple}Description} \\
\midrule
\texttt{ECALL}  & Environment Call: requests a service from the operating system or execution environment (e.g., a system call). The specific service is determined by the values in \texttt{a0}--\texttt{a7}. \\[4pt]
\texttt{EBREAK} & Environment Breakpoint: transfers control to a debugger. Used to set software breakpoints during debugging. \\
\bottomrule
\end{tabular}
\end{center}

\begin{keybox}
\textbf{\color{gold}ECALL vs.\ EBREAK:}
These instructions do not behave like normal arithmetic or memory operations. Instead, they hand control to something \textbf{outside the ordinary instruction flow}.

\begin{itemize}
  \item \texttt{ECALL} is used when the program wants the execution environment to do something on its behalf.
  \item \texttt{EBREAK} is used when the program wants to stop under debugger control.
\end{itemize}

The exact behavior of \texttt{ECALL} depends on the simulator, operating system, or test environment being used, while \texttt{EBREAK} is mainly a debugging aid.
\end{keybox}

\newpage

% =====================================================
% CHAPTER 3: COMPLETE ISA REFERENCE TABLE
% =====================================================
\chapter{Complete RV32I Instruction Reference}

This chapter provides the \textbf{full RV32I instruction encoding table} as specified in the official RISC-V ISA manual. Every instruction is listed with its type, opcode, funct3, funct7 (where applicable), and a concise description.

\begin{center}
  \includegraphics[width=0.98\textwidth]{images/isa_map.png}
\end{center}

\section{Master Encoding Table}

\begin{center}
\renewcommand{\arraystretch}{1.18}
\footnotesize
\setlength{\LTpre}{6pt}
\setlength{\LTpost}{6pt}
\begin{longtable}{
  >{\ttfamily}l  % Instruction
  c              % Type
  >{\ttfamily}c  % opcode
  >{\ttfamily}c  % funct3
  >{\ttfamily}c  % funct7/imm
  p{5.5cm}       % Description
}
\toprule
\normalfont\textbf{\color{purple}Instruction} &
\textbf{\color{purple}Type} &
\textbf{\color{purple}opcode} &
\textbf{\color{purple}funct3} &
\textbf{\color{purple}funct7/imm} &
\textbf{\color{purple}Description} \\
\midrule
\endfirsthead

\toprule
\normalfont\textbf{\color{purple}Instruction} &
\textbf{\color{purple}Type} &
\textbf{\color{purple}opcode} &
\textbf{\color{purple}funct3} &
\textbf{\color{purple}funct7/imm} &
\textbf{\color{purple}Description} \\
\midrule
\endhead

\midrule
\multicolumn{6}{r}{\footnotesize\itshape Continued on next page\ldots} \\
\bottomrule
\endfoot

\bottomrule
\endlastfoot

%% ─── R-TYPE ──────────────────────────────────────────
\rowcolor{funct7bg!30}
\multicolumn{6}{l}{\textbf{\color{rtypecolor}\normalfont R-Type --- Register-Register Operations (opcode 0110011)}} \\
ADD  rd,rs1,rs2   & R & 0110011 & 000 & 0000000 & Add: rd = rs1 + rs2 \\
SUB  rd,rs1,rs2   & R & 0110011 & 000 & 0100000 & Subtract: rd = rs1 $-$ rs2 \\
SLL  rd,rs1,rs2   & R & 0110011 & 001 & 0000000 & Shift left logical by rs2[4:0] \\
SLT  rd,rs1,rs2   & R & 0110011 & 010 & 0000000 & Set if less than (signed) \\
SLTU rd,rs1,rs2   & R & 0110011 & 011 & 0000000 & Set if less than (unsigned) \\
XOR  rd,rs1,rs2   & R & 0110011 & 100 & 0000000 & Bitwise exclusive OR \\
SRL  rd,rs1,rs2   & R & 0110011 & 101 & 0000000 & Shift right logical by rs2[4:0] \\
SRA  rd,rs1,rs2   & R & 0110011 & 101 & 0100000 & Shift right arithmetic by rs2[4:0] \\
OR   rd,rs1,rs2   & R & 0110011 & 110 & 0000000 & Bitwise OR \\
AND  rd,rs1,rs2   & R & 0110011 & 111 & 0000000 & Bitwise AND \\
\midrule

%% ─── I-TYPE ARITHMETIC ──────────────────────────────
\rowcolor{rs1bg!30}
\multicolumn{6}{l}{\textbf{\color{itypecolor}\normalfont I-Type --- Immediate Arithmetic (opcode 0010011)}} \\
ADDI  rd,rs1,imm  & I & 0010011 & 000 & ---     & Add immediate \\
SLTI  rd,rs1,imm  & I & 0010011 & 010 & ---     & Set if $<$ imm (signed) \\
SLTIU rd,rs1,imm  & I & 0010011 & 011 & ---     & Set if $<$ imm (unsigned) \\
XORI  rd,rs1,imm  & I & 0010011 & 100 & ---     & XOR with immediate \\
ORI   rd,rs1,imm  & I & 0010011 & 110 & ---     & OR with immediate \\
ANDI  rd,rs1,imm  & I & 0010011 & 111 & ---     & AND with immediate \\
SLLI  rd,rs1,shamt& I & 0010011 & 001 & 0000000 & Shift left logical by shamt \\
SRLI  rd,rs1,shamt& I & 0010011 & 101 & 0000000 & Shift right logical by shamt \\
SRAI  rd,rs1,shamt& I & 0010011 & 101 & 0100000 & Shift right arithmetic by shamt \\
\midrule

%% ─── I-TYPE LOADS ───────────────────────────────────
\rowcolor{rs2bg!30}
\multicolumn{6}{l}{\textbf{\color{stypecolor}\normalfont I-Type --- Load Instructions (opcode 0000011)}} \\
LB   rd,imm(rs1)  & I & 0000011 & 000 & ---     & Load byte, sign-extend \\
LH   rd,imm(rs1)  & I & 0000011 & 001 & ---     & Load halfword, sign-extend \\
LW   rd,imm(rs1)  & I & 0000011 & 010 & ---     & Load word (32-bit) \\
LBU  rd,imm(rs1)  & I & 0000011 & 100 & ---     & Load byte, zero-extend \\
LHU  rd,imm(rs1)  & I & 0000011 & 101 & ---     & Load halfword, zero-extend \\
\midrule

%% ─── S-TYPE STORES ──────────────────────────────────
\rowcolor{rs2bg!30}
\multicolumn{6}{l}{\textbf{\color{stypecolor}\normalfont S-Type --- Store Instructions (opcode 0100011)}} \\
SB   rs2,imm(rs1) & S & 0100011 & 000 & ---     & Store byte (8-bit) \\
SH   rs2,imm(rs1) & S & 0100011 & 001 & ---     & Store halfword (16-bit) \\
SW   rs2,imm(rs1) & S & 0100011 & 010 & ---     & Store word (32-bit) \\
\midrule

%% ─── B-TYPE BRANCHES ────────────────────────────────
\rowcolor{funct3bg!30}
\multicolumn{6}{l}{\textbf{\color{btypecolor}\normalfont B-Type --- Conditional Branches (opcode 1100011)}} \\
BEQ  rs1,rs2,off  & B & 1100011 & 000 & ---     & Branch if rs1 $=$ rs2 \\
BNE  rs1,rs2,off  & B & 1100011 & 001 & ---     & Branch if rs1 $\neq$ rs2 \\
BLT  rs1,rs2,off  & B & 1100011 & 100 & ---     & Branch if rs1 $<$ rs2 (signed) \\
BGE  rs1,rs2,off  & B & 1100011 & 101 & ---     & Branch if rs1 $\geq$ rs2 (signed) \\
BLTU rs1,rs2,off  & B & 1100011 & 110 & ---     & Branch if rs1 $<$ rs2 (unsigned) \\
BGEU rs1,rs2,off  & B & 1100011 & 111 & ---     & Branch if rs1 $\geq$ rs2 (unsigned) \\
\midrule

%% ─── U-TYPE ─────────────────────────────────────────
\rowcolor{rdbg!30}
\multicolumn{6}{l}{\textbf{\color{utypecolor}\normalfont U-Type --- Upper Immediate}} \\
LUI   rd,imm      & U & 0110111 & --- & ---     & Load upper immediate (bits [31:12]) \\
AUIPC rd,imm      & U & 0010111 & --- & ---     & Add upper immediate to PC \\
\midrule

%% ─── J-TYPE / JALR ──────────────────────────────────
\rowcolor{opcodebg!30}
\multicolumn{6}{l}{\textbf{\color{jtypecolor}\normalfont J-Type / I-Type --- Jump Instructions}} \\
JAL  rd,offset     & J & 1101111 & --- & ---     & Jump and link (PC-relative) \\
JALR rd,rs1,off    & I & 1100111 & 000 & ---     & Jump and link register (indirect) \\
\midrule

%% ─── SYSTEM ─────────────────────────────────────────
\rowcolor{gold!15}
\multicolumn{6}{l}{\textbf{\color{gold!80!black}\normalfont System Instructions (opcode 1110011)}} \\
ECALL              & I & 1110011 & 000 & 000000000000 & Environment call (syscall) \\
EBREAK             & I & 1110011 & 000 & 000000000001 & Environment breakpoint \\

\end{longtable}
\end{center}

\newpage

\section{Instruction Count Summary}

\begin{center}
\renewcommand{\arraystretch}{1.3}
\begin{tabular}{l c l}
\toprule
\textbf{\color{purple}Category} & \textbf{\color{purple}Count} & \textbf{\color{purple}Instructions} \\
\midrule
R-Type (Register)        & 10 & ADD, SUB, SLL, SLT, SLTU, XOR, SRL, SRA, OR, AND \\
I-Type (Imm.\ Arith.)   & 9  & ADDI, SLTI, SLTIU, XORI, ORI, ANDI, SLLI, SRLI, SRAI \\
I-Type (Load)            & 5  & LB, LH, LW, LBU, LHU \\
S-Type (Store)           & 3  & SB, SH, SW \\
B-Type (Branch)          & 6  & BEQ, BNE, BLT, BGE, BLTU, BGEU \\
U-Type (Upper Imm.)      & 2  & LUI, AUIPC \\
J-Type (JAL)             & 1  & JAL \\
I-Type (JALR)            & 1  & JALR \\
System                   & 2  & ECALL, EBREAK \\
\midrule
\textbf{Total}           & \textbf{39} & \\
\bottomrule
\end{tabular}
\end{center}

\vspace{0.5cm}

\begin{keybox}
\textbf{\color{gold}Design Implication:}
These 39 instructions (often quoted as ``40 instructions'' when including FENCE) represent the \textit{entire} base RV32I ISA. From a hardware designer's perspective, each instruction translates to a specific combination of control signals. Understanding every instruction's encoding and semantics is the foundation for designing the datapath, control unit, and decoder of the YARP processor.
\end{keybox}

\newpage

% =====================================================
% CHAPTER 4: INSTRUCTION DETAILS & EXAMPLES
% =====================================================
\chapter{Instruction Details \& Examples}

This chapter provides worked examples showing how each instruction type is encoded as a 32-bit machine word and what it does conceptually.

\section{R-Type Example: ADD x5, x6, x7}

\begin{center}
\renewcommand{\arraystretch}{1.5}
\begin{tabular}{|>{\centering\arraybackslash\columncolor{funct7bg}}p{2.2cm}|>{\centering\arraybackslash\columncolor{rs2bg}}p{1.6cm}|>{\centering\arraybackslash\columncolor{rs1bg}}p{1.6cm}|>{\centering\arraybackslash\columncolor{funct3bg}}p{1.5cm}|>{\centering\arraybackslash\columncolor{rdbg}}p{1.6cm}|>{\centering\arraybackslash\columncolor{opcodebg}}p{2.2cm}|}
\hline
\texttt{0000000} & \texttt{00111} & \texttt{00110} & \texttt{000} & \texttt{00101} & \texttt{0110011} \\
\hline
funct7 & \texttt{x7} & \texttt{x6} & ADD & \texttt{x5} & OP \\
\hline
\end{tabular}
\end{center}

\textbf{Binary:} \texttt{0000000\,00111\,00110\,000\,00101\,0110011}

\textbf{Hex:} \texttt{0x007302B3}

\textbf{Semantics:} $\texttt{x5} \leftarrow \texttt{x6} + \texttt{x7}$. Reads registers \texttt{x6} and \texttt{x7}, computes their sum, and writes the result to \texttt{x5}.

\section{I-Type Example: ADDI x10, x0, 42}

\begin{center}
\renewcommand{\arraystretch}{1.5}
\begin{tabular}{|>{\centering\arraybackslash\columncolor{immbg}}p{3.8cm}|>{\centering\arraybackslash\columncolor{rs1bg}}p{1.6cm}|>{\centering\arraybackslash\columncolor{funct3bg}}p{1.5cm}|>{\centering\arraybackslash\columncolor{rdbg}}p{1.6cm}|>{\centering\arraybackslash\columncolor{opcodebg}}p{2.2cm}|}
\hline
\texttt{000000101010} & \texttt{00000} & \texttt{000} & \texttt{01010} & \texttt{0010011} \\
\hline
imm = 42 & \texttt{x0} & ADDI & \texttt{x10} & OP-IMM \\
\hline
\end{tabular}
\end{center}

\textbf{Semantics:} $\texttt{x10} \leftarrow \texttt{x0} + 42 = 0 + 42 = 42$. This is the standard way to load a small constant into a register.

\section{S-Type Example: SW x15, 8(x2)}

Store word: writes the contents of \texttt{x15} to memory at address $\texttt{x2} + 8$.

\begin{center}
\renewcommand{\arraystretch}{1.5}
\begin{tabular}{|>{\centering\arraybackslash\columncolor{immbg}}p{2.2cm}|>{\centering\arraybackslash\columncolor{rs2bg}}p{1.6cm}|>{\centering\arraybackslash\columncolor{rs1bg}}p{1.6cm}|>{\centering\arraybackslash\columncolor{funct3bg}}p{1.5cm}|>{\centering\arraybackslash\columncolor{immbg}}p{1.6cm}|>{\centering\arraybackslash\columncolor{opcodebg}}p{2.2cm}|}
\hline
\texttt{0000000} & \texttt{01111} & \texttt{00010} & \texttt{010} & \texttt{01000} & \texttt{0100011} \\
\hline
imm[11:5]=0 & \texttt{x15} & \texttt{x2} & SW & imm[4:0]=8 & STORE \\
\hline
\end{tabular}
\end{center}

\textbf{Semantics:} $\text{Mem}[\texttt{x2}+8] \leftarrow \texttt{x15}$

\section{B-Type Example: BEQ x1, x2, 16}

Branch to $\text{PC}+16$ if \texttt{x1} equals \texttt{x2}.

\textbf{Offset = 16:} binary \texttt{0\,0000001000\,0} $\rightarrow$ imm[12:1] = \texttt{0\_000000\_1000\_0}

The scattered immediate fields encode: imm[12]=0, imm[11]=0, imm[10:5]=000000, imm[4:1]=1000.

\textbf{Semantics:} If $\texttt{x1} = \texttt{x2}$, then $\text{PC} \leftarrow \text{PC} + 16$; else $\text{PC} \leftarrow \text{PC} + 4$.

\section{U-Type Example: LUI x5, 0xDEADC}

\textbf{Semantics:} $\texttt{x5} \leftarrow \texttt{0xDEADC} \ll 12 = \texttt{0xDEADC000}$. Places the upper 20 bits directly.

\section{J-Type Example: JAL ra, 100}

\textbf{Semantics:}
\begin{enumerate}
  \item $\texttt{ra} \leftarrow \text{PC} + 4$ (save return address)
  \item $\text{PC} \leftarrow \text{PC} + 100$ (jump to target)
\end{enumerate}

This is a function call: the processor remembers where to return (in \texttt{ra}) and jumps to the function body that is 100 bytes away.

\section{Common Pseudoinstructions}

RISC-V assemblers recognize \textbf{pseudoinstructions} --- mnemonics that don't correspond to a single hardware instruction but are expanded into one or more real instructions:

\begin{center}
\renewcommand{\arraystretch}{1.25}
\begin{tabular}{l l p{5.5cm}}
\toprule
\textbf{\color{purple}Pseudo} & \textbf{\color{purple}Expansion} & \textbf{\color{purple}Meaning} \\
\midrule
\texttt{nop}          & \texttt{addi x0, x0, 0}    & No operation \\
\texttt{li rd, imm}   & \texttt{lui + addi}         & Load 32-bit immediate \\
\texttt{mv rd, rs}    & \texttt{addi rd, rs, 0}     & Copy register \\
\texttt{not rd, rs}   & \texttt{xori rd, rs, -1}    & Bitwise NOT \\
\texttt{neg rd, rs}   & \texttt{sub rd, x0, rs}     & Two's complement negate \\
\texttt{j offset}     & \texttt{jal x0, offset}     & Unconditional jump (no link) \\
\texttt{ret}          & \texttt{jalr x0, ra, 0}     & Return from function \\
\texttt{call label}   & \texttt{auipc + jalr}       & Far function call \\
\texttt{beqz rs, off} & \texttt{beq rs, x0, off}    & Branch if zero \\
\texttt{bnez rs, off} & \texttt{bne rs, x0, off}    & Branch if not zero \\
\texttt{bgt rs,rt,off}& \texttt{blt rt, rs, off}    & Branch if greater than \\
\texttt{ble rs,rt,off}& \texttt{bge rt, rs, off}    & Branch if less or equal \\
\bottomrule
\end{tabular}
\end{center}

\begin{infobox}
\textbf{\color{accentblue}Why pseudoinstructions matter:}
Pseudoinstructions make assembly code easier to read and write, while keeping the hardware simple. The assembler transparently expands them. For verification, only the \textbf{real} (base) instructions need to be implemented in hardware.
\end{infobox}

\end{document}
